

We consider fitting this model by aggregating the data at each time point and treating the averages as a time series with covariates. 100 simulations were conducted with the following parameters:

$$ \beta = (43, -.001, .25), \qquad \sigma_\varepsilon^2 = 100, \qquad \phi = \theta = .5, \qquad \sigma_\eta^2 = 1  $$

Some sample rows of simulated data are provided below:


\input{images/output/arma_cov_head.txt}

After aggregating at each time step, we obtain the following condensed dataset, from which a time series with covariates is then fit:

\input{images/output/arma_cov_head_daily.txt}

Note that the \textit{x} column of each time step sample has a value close to $0.5$, with some variability. If we then fit a time series with covariates to the condensed dataset, we obtain results with empirical means displayed in the table below. The ``Significance Frequency" column displays the fraction of the hundred simulations in which the corresponding parameter estimate was statistically significant at the $\alpha = .05$ level.


\input{images/output/arma_cov_df.txt}

Note that the effect of the dummy indicator variable is deemed significant only $11\%$ of the time. This example highlights the shortcomings of fitting a time series with covariates. Also of note is how skewed the sample mean of $\hat{\theta}$ is, in particular how close to zero it is. We will see later why this is actually expected. For now, we use these results to motivate the need for a different approach. With that need in mind, we next review some of the key developments in modeling repeated cross-sectional data.